% CSE 715 — Computer Graphics
% Octrees / Quadtrees / BSP-Trees / Fractal-Geometry Methods
% Source: Computer Graphics C Version (Hearn & Baker), Ch10 §10-16 (Octrees), §10-17 (BSP Trees), §10-18 (Fractal-Geometry Methods)
% University of Chittagong, 7th Semester B.Sc. Engineering
%
% Compile with: pdflatex OctreeFractal_Solutions.tex (run twice for TOC)

\documentclass[12pt, a4paper]{article}

\usepackage[left=2.5cm, right=2.5cm, top=2.8cm, bottom=2.8cm, headheight=15pt]{geometry}
\usepackage{amsmath, amssymb, mathtools}
\usepackage{tikz}
\usetikzlibrary{arrows.meta, positioning, shapes.geometric, calc, backgrounds, fit}
\usepackage{booktabs}
\usepackage{array}
\usepackage{xcolor}
\usepackage{enumitem}
\usepackage{fancyhdr}
\usepackage{titlesec}
\usepackage{mdframed}
\usepackage{tabularx}
\usepackage{hyperref}

\definecolor{sectionblue}{RGB}{10,60,110}
\definecolor{boxbg}{RGB}{242,247,255}
\definecolor{answerbg}{RGB}{240,255,240}
\definecolor{notesbg}{RGB}{255,250,230}

\mdfdefinestyle{solutionframe}{linecolor=sectionblue, backgroundcolor=boxbg, roundcorner=4pt, linewidth=1.4pt, innertopmargin=7pt, innerbottommargin=7pt, innerleftmargin=9pt, innerrightmargin=9pt}
\mdfdefinestyle{answerframe}{linecolor=green!55!black, backgroundcolor=answerbg, roundcorner=4pt, linewidth=1.4pt, innertopmargin=7pt, innerbottommargin=7pt, innerleftmargin=9pt, innerrightmargin=9pt}
\mdfdefinestyle{notesframe}{linecolor=orange!70!black, backgroundcolor=notesbg, roundcorner=4pt, linewidth=1pt, innertopmargin=5pt, innerbottommargin=5pt, innerleftmargin=8pt, innerrightmargin=8pt}
\newenvironment{solutionbox}{\begin{mdframed}[style=solutionframe]}{\end{mdframed}}
\newenvironment{answerbox}{\begin{mdframed}[style=answerframe]}{\end{mdframed}}
\newenvironment{notesbox}{\begin{mdframed}[style=notesframe]}{\end{mdframed}}

\pagestyle{fancy}
\fancyhf{}
\lhead{\small\color{sectionblue} CSE 715 --- Octrees/BSP-Trees/Fractal Geometry}
\rhead{\small\color{sectionblue} Syllabus-Mandated Coverage (Hearn \& Baker)}
\cfoot{\small\thepage}
\renewcommand{\headrulewidth}{0.5pt}

\titleformat{\section}{\Large\bfseries\color{sectionblue}}{}{0em}{}[\vspace{-2pt}\color{sectionblue}\rule{\linewidth}{0.8pt}]
\titleformat{\subsection}{\large\bfseries\color{sectionblue!85!black}}{}{0em}{}

\begin{document}

\begin{center}
  {\LARGE\bfseries\color{sectionblue} CSE 715: Computer Graphics}\\[0.6em]
  {\Large\bfseries Octrees, Quadtrees, BSP-Trees \& Fractal-Geometry Methods}\\[0.3em]
  {\large Syllabus-Mandated Coverage --- ``the 5 no-resource topics''}\\[0.2em]
  {\normalsize University of Chittagong \quad|\quad 7\textsuperscript{th} Semester B.Sc.\ (Engg.) Examination}\\[0.3em]
  {\small Source: \emph{Computer Graphics C Version} (Hearn \& Baker), Chapter 10 --- \S10-16 (Octrees, incl.\ Quadtrees), \S10-17 (BSP Trees), \S10-18 (Fractal-Geometry Methods)}
\end{center}

\vspace{0.3em}\noindent\rule{\linewidth}{1pt}\vspace{0.5em}

\begin{notesbox}
\textbf{Why this document looks different from the other chapter solutions:} this topic has \textbf{0/5 hits} in the 2020--2024 past papers --- there is no PYQ history to analyze. \texttt{Syllabus.txt} explicitly carves out an exception: \emph{``from the book: Computer Graphics C version, only the following topics: 1. Octrees (how it works) page 359, Quadtree, BST-Tree 10-17, Fractal Geometry 10-18, B-spline Curve, types of Shading.''} (B-spline and types-of-Shading are already covered in [[wiki/bezier-bspline-hermite|Ch9]] and [[wiki/color-shading-models|Ch11]] respectively --- this document covers the remaining, previously-uncovered piece: Octrees/BSP-Trees/Fractal Geometry.) Given the exam is tomorrow and this is the \textbf{lowest-yield item in the entire 51-topic tracker}, treat this as a stretch-goal skim only --- every other TIER1/TIER2 topic and, above all, timed past-paper practice come first.
\end{notesbox}

\tableofcontents
\newpage

%=====================================================================
\section{Octrees (\S10-16)}
%=====================================================================
\begin{answerbox}
An \textbf{octree} is a hierarchical tree structure used to represent solid 3D objects by recursively subdividing a bounding region of space into \textbf{octants} (eight equal sub-cubes). Each node has \textbf{8 data elements}, one per octant; the smallest unit of 3D space is a \textbf{voxel} (volume element).
\begin{itemize}[noitemsep]
  \item If all voxels in an octant are the same type (\textbf{homogeneous}), that octant's data element simply stores the type value (a leaf).
  \item If an octant is \textbf{heterogeneous}, it is subdivided further into 8 sub-octants, and the data element instead stores a pointer to the next tree node.
  \item Subdivision continues until every octant is homogeneous or a resolution limit is reached. Each node has 0 to 8 immediate descendants.
  \item \textbf{Quadtree} is the 2D analogue: a region (usually square) is divided into \textbf{4 quadrants} per node instead of 8 octants; used for 2D image/region encoding and as a stepping stone for explaining octrees.
\end{itemize}
\textbf{How it works (generation):} define a bounding box around the object using its min/max coordinates; test the box octant-by-octant (from any input representation: polygon mesh, curved patches, or CSG solids); homogeneous octants become leaves, heterogeneous ones recurse.
\textbf{Uses:} set operations (union/intersection/difference) on two octree objects by scanning/comparing octants; 3D rotations by transforming occupied octants; visible-surface identification by searching octants front-to-back (first hit = visible) and transferring the result to a quadtree for display; medical imaging and other applications needing object cross-sections; reduces storage via spatial coherence (large homogeneous regions collapse to one node).
\end{answerbox}

%=====================================================================
\section{BSP Trees (\S10-17)}
%=====================================================================
\begin{answerbox}
A \textbf{binary space-partitioning (BSP) tree} subdivides a scene into \textbf{two} partitions at each step (vs.\ octrees' fixed eight), using a cutting \textbf{plane that can be at any position and orientation} — unlike an octree, whose three cutting planes at each step are always mutually perpendicular and aligned with the coordinate axes.
\begin{itemize}[noitemsep]
  \item \textbf{Adaptive:} since the cutting plane's position/orientation can be chosen to fit the actual spatial distribution of objects, a BSP tree can be \textbf{shallower} than the equivalent octree for the same scene, reducing tree-search time.
  \item \textbf{Uses:} identifying visible surfaces, and space-partitioning for \textbf{ray-tracing algorithms} (accelerating ray/object intersection tests by quickly discarding whole partitions a ray cannot reach) — a direct link to [[wiki/ray-tracing|Ch12's]] spatial-subdivision efficiency technique.
\end{itemize}
\end{answerbox}

%=====================================================================
\section{Fractal-Geometry Methods (\S10-18)}
%=====================================================================

\subsection{Motivation and Core Definition}
\begin{answerbox}
Euclidean-geometry methods (equations) describe manufactured objects with smooth surfaces/regular shapes well, but cannot realistically model \textbf{natural} objects (mountains, clouds, coastlines) that have irregular, fragmented features. \textbf{Fractal-geometry methods} instead describe objects \textbf{procedurally} (a repeated transformation/generation procedure) rather than with equations.
\textbf{A fractal object has two characteristics:}
\begin{enumerate}[noitemsep]
  \item \textbf{Infinite detail at every point} — zooming into a fractal (e.g.\ a mountain outline) keeps revealing new jagged detail, unlike a Euclidean shape which eventually looks smooth on zooming in.
  \item \textbf{Self-similarity} between the object's parts and its overall features — the exact form of self-similarity depends on which fractal representation is chosen.
\end{enumerate}
Since true infinite detail can't be displayed, graphics implementations apply the generating transformation a \textbf{finite} number of times (enough to exceed pixel resolution).
\end{answerbox}

\subsection{Fractal-Generation Procedure}
\begin{answerbox}
Given an initial point (or primitive) $P_0$, repeatedly apply a transformation function $F$:
$$P_1 = F(P_0), \quad P_2 = F(P_1), \quad P_3 = F(P_2), \ \ldots$$
$F$ may be geometric (scaling/translation/rotation), a nonlinear coordinate transform, deterministic, or randomized at each iteration. More iterations $\Rightarrow$ more detail, up to the display's pixel resolution.
\end{answerbox}

\subsection{Classification of Fractals}
\begin{answerbox}
\begin{center}
\begin{tabular}{p{3.2cm}p{10.2cm}}
\toprule
\textbf{Class} & \textbf{Defining property} \\
\midrule
Self-similar & Parts are scaled-down copies of the whole, using a single scaling factor $s$ for all subparts (or different factors per subpart) \\
Statistically self-similar & Self-similar, plus random variation applied to each scaled-down subpart (models trees, shrubs, plants) \\
Self-affine & Parts scaled with \emph{different} factors $s_x,s_y,s_z$ per coordinate direction \\
Statistically self-affine & Self-affine + random variation (models terrain, water, clouds) \\
Invariant fractal sets & Formed with nonlinear transformations: \textbf{self-squaring} (e.g.\ the Mandelbrot set, from squaring functions in complex space) and \textbf{self-inverse} (from inversion procedures) \\
\bottomrule
\end{tabular}
\end{center}
\end{answerbox}

\subsection{Fractal Dimension}
\begin{answerbox}
The \textbf{fractal dimension} $D$ (a.k.a.\ \emph{fractional dimension} — the source of the name ``fractal'') measures the roughness/fragmentation of an object; more jagged objects have larger $D$. For a self-similar fractal built with a single scaling factor $s$, subdividing into $n$ subparts satisfies $ns^D=1$, giving
$$D = \frac{\ln n}{\ln(1/s)}$$
For different scaling factors $s_k$ per subpart ($k=1,\dots,n$), $D$ solves the implicit relation $\sum_{k=1}^n s_k^D = 1$.
\textbf{Worked examples from the text:}
\begin{itemize}[noitemsep]
  \item \textbf{Koch curve/snowflake:} each segment $\to$ 4 segments, scaling factor $s=1/3$: $D=\dfrac{\ln4}{\ln3}\approx1.2619$.
  \item \textbf{Fractal tetrahedron surface:} each face $\to$ 6 smaller faces, scaling factor $s=1/2$: $D=\dfrac{\ln6}{\ln2}\approx2.58496$ (a fairly fragmented surface).
\end{itemize}
\textbf{Dimension ranges (Euclidean dimension is always exceeded):} a fractal \emph{curve} has $D>1$ ($D=2$: a Peano curve completely fills a 2D region; $2<D<3$: self-intersecting, can cover an infinite area). A fractal \emph{surface} has $2<D\le3$ ($D=3$: fills a volume; $D>3$: overlapping coverage). A fractal \emph{solid} has $3<D\le4$.
\textbf{Estimating $D$ for irregular/non-simple shapes:} \emph{box-covering methods} — cover the object with $n$ small boxes (squares/cubes) of a given scaling factor and use $n,s$ in the formula above; the more general version is the \textbf{Hausdorff--Besicovitch dimension}.
\end{answerbox}

\subsection*{Illustrative Exercise --- Fractal Dimension of a Custom Generator}
\begin{notesbox}
Illustrative exercise (not from a past paper), same style as the worked Koch-curve example, to confirm the formula $D=\ln n/\ln(1/s)$.
\end{notesbox}
\begin{solutionbox}
\textbf{Given} a self-similar generator that replaces each straight segment with \textbf{5} equal-length segments, each scaled by $s=1/5$, find the fractal dimension.
\end{solutionbox}
\begin{answerbox}
Using $D=\dfrac{\ln n}{\ln(1/s)}$ with $n=5,\ s=1/5$: $D=\dfrac{\ln5}{\ln5}=1$. This makes sense: replacing 1 segment with 5 collinear-scale segments covering the same total length end-to-end produces no extra jaggedness beyond a straight line (unlike the Koch curve's bump, which pushes 2 of the 4 sub-segments off the original line) — a sanity check that $D=1$ corresponds to an (effectively) non-fractal, ordinary Euclidean line.
\end{answerbox}

%=====================================================================
\section{Quick Summary}
%=====================================================================
\begin{center}
\renewcommand{\arraystretch}{1.3}
\begin{tabular}{p{3.4cm}|p{10cm}}
\toprule
\textbf{Concept} & \textbf{One-line memory hook} \\
\midrule
Octree & 3D space split into \textbf{8} octants/node, axis-aligned; homogeneous octants are leaves \\
Quadtree & 2D analogue of octree --- \textbf{4} quadrants/node \\
BSP tree & \textbf{2} partitions/step, cutting plane can be \emph{any} position/orientation --- adaptive, shallower than octree; used in ray-tracing space subdivision \\
Fractal object & Infinite detail + self-similarity; built by repeating a transformation $F$ \\
Fractal dimension & $D=\ln n/\ln(1/s)$ --- always exceeds the Euclidean dimension; higher $D$ = more jagged/fragmented \\
\bottomrule
\end{tabular}
\end{center}
\begin{notesbox}
\textbf{Exam-day reality check:} this topic has never appeared in 5 years of past papers and is the lowest-yield item in the entire tracker. If time is short, this is the first thing to skip back from --- Z-Buffer, Cohen-Sutherland, and the other TIER1/TIER2 topics (and, most of all, timed past-paper practice) come first.
\end{notesbox}

\end{document}
